
\mxunit{stmg-c02}{Fonctions exponentielles}
\mxskills{modéliser une évolution continue, changer de registre, calculer, interpréter}
\section{Cours de référence}
Pour $a>0$, la fonction $x\mapsto a^x$ prolonge à tous les réels les puissances entières de $a$. Elle vérifie $a^0=1$ et $a^{-x}=\dfrac{1}{a^x}$.

\begin{center}
\begin{tikzpicture}
\begin{axis}[width=.78\textwidth,height=6cm,axis lines=middle,xmin=-3,xmax=4,ymin=-2,ymax=6,grid=major,xlabel={$x$},ylabel={$y$},legend style={at={(.03,.97)},anchor=north west,draw=none,fill=white}]
\addplot[mxteacher,very thick,samples=100,domain=-3:4]{1.6^x};\addlegendentry{$\num{1,6}^x$}
\addplot[mxwarn,very thick,samples=100,domain=-3:4]{0.65^x};\addlegendentry{$\num{0,65}^x$}
\addplot[mxmint,very thick,samples=100,domain=-3:4]{-1.2*1.35^x};\addlegendentry{$-\num{1,2}\times\num{1,35}^x$}
\end{axis}
\end{tikzpicture}
\end{center}

\begin{keypoint}
Si $a>1$, $x\mapsto a^x$ est strictement croissante. Si $0<a<1$, elle est strictement décroissante. Pour $f(x)=ka^x$, multiplier par $k<0$ inverse le sens de variation. Dans tous les cas, $a^x>0$.
\end{keypoint}

Pour $a>0$ et tous réels $x,y$ :
\[
a^{x+y}=a^xa^y,\qquad a^{x-y}=\frac{a^x}{a^y},\qquad (a^x)^y=a^{xy}.
\]

\begin{pitfall}
En général, $a^x+a^y\neq a^{x+y}$. Les propriétés portent sur les produits et les quotients, pas sur les sommes.
\end{pitfall}

\subsection{Évolutions et taux moyen}
Une quantité de valeur initiale $V_0$ soumise à une évolution relative constante de coefficient $q>0$ se modélise par $V(t)=V_0q^t$. Si un coefficient global $Q>0$ résulte de $n$ évolutions identiques, le coefficient moyen est $q=Q^{\frac1n}$ et le taux moyen est $\tau=q-1$.

\begin{example}
Un chiffre d'affaires passe de $\num{80000}$ à $\num{95000}$ euros en quatre ans. Le coefficient global vaut $Q=\dfrac{\num{95000}}{\num{80000}}=\num{1,1875}$. Le coefficient annuel moyen est $q=\num{1,1875}^{\frac14}\approx\num{1,0439}$ : le taux annuel moyen est donc environ $\pc{4,39}$.
\end{example}

\subsection{Résoudre et interpréter}
Lorsque les deux membres peuvent être écrits avec la même base, on utilise l'injectivité : pour $a>0$ et $a\neq1$, $a^u=a^v$ équivaut à $u=v$. Sinon, une résolution graphique ou numérique donne une valeur approchée ; le logarithme décimal, étudié ensuite, donnera une formule.

\mxsession{7}{\mxdate*{2026-10-06}}{B : 2 h -- C : 1 h}{Définition et sens de variation de $a^x$}{Fonctions exponentielles}
\begin{objectives}\begin{itemize}\item Donner du sens à $a^x$.\item Relier base et variation.\item Lire des courbes.\end{itemize}\end{objectives}
\begin{fluency}[6 à 8 minutes, sans calculatrice]\begin{enumerate}\item Écrire le coefficient multiplicateur d'une hausse de $\pc{10}$. \mxblank{$1+\dfrac{10}{100}=\num{1,10}$}\item Comparer $\num{1,04}^3$ et $\num{1,04}^7$. \mxblank{$\num{1,04}^3<\num{1,04}^7$}\item Calculer $2^4\times2^3$. \mxblank{$2^7$}\item Écrire $5^{-2}$ sous forme de fraction. \mxblank{$\dfrac1{25}$}\item La fonction $x\mapsto\num{0,7}^x$ est-elle croissante ? \mxblank{Non, elle est décroissante.}\end{enumerate}\end{fluency}
\begin{prerequisites}Relire la dernière trace écrite et savoir effectuer les calculs élémentaires mobilisés aujourd'hui. Avant l'activité, faire reformuler les mots-clés : donner du sens à $a^x$ et relier base et variation.\end{prerequisites}

\begin{activity}
À l'aide d'une calculatrice, compléter un tableau de valeurs de $\num{1,5}^x$ et $\num{0,7}^x$ pour $x=-2,-1,0,1,2,3$. Tracer les points, conjecturer les variations et identifier le point commun aux deux courbes.
\begin{center}\begin{tikzpicture}\begin{axis}[width=.68\textwidth,height=5cm,axis lines=middle,xmin=-3,xmax=4,ymin=0,ymax=5,grid=major]\addplot[mxteacher,very thick,samples=100,domain=-3:4]{1.5^x};\addplot[mxwarn,very thick,samples=100,domain=-3:4]{.7^x};\node[mxteacher,anchor=south east] at (axis cs:3.2,3.9) {$\mathcal C_1$};\node[mxwarn,anchor=south west] at (axis cs:-2.9,3.3) {$\mathcal C_2$};\end{axis}\end{tikzpicture}\end{center}
\end{activity}
\begin{solution}$a^0=1$, donc les courbes passent par $\pt{0;1}$. La première croît et la seconde décroît. Faire verbaliser la donnée qui déclenche le choix de méthode.\end{solution}

\begin{summary}Pour $a>0$, $a^x$ est défini pour tout réel $x$. La fonction est croissante si $a>1$ et décroissante si $0<a<1$.\end{summary}

\begin{example}
Reprendre les données de l'activité et rédiger la solution dans l'ordre : définition des objets, propriété utilisée, calcul exact ou approché, puis phrase de réponse. La réponse modèle figure dans l'édition enseignant.
\end{example}
\begin{solution}$a^0=1$, donc les courbes passent par $\pt{0;1}$. La première croît et la seconde décroît.\end{solution}

\begin{exercise}[essentiel]
\begin{enumerate}\item Simplifier $3^{x+4}\div3^x$ et $(2^x)^3$.\item Étudier le sens de variation de $f(x)=4\times\num{1,04}^x$ et $g(x)=-4\times\num{0,8}^x$.\item Une valeur initiale de $1350$ augmente de $\pc{4}$ par an. Donner le modèle et calculer la valeur après $6$ ans.\end{enumerate}
\end{exercise}
\begin{solution}$3^4$ et $2^{3x}$. $f$ est croissante car $\num{1,04}>1$ et $4>0$ ; $\num{0,8}^x$ est décroissante et le coefficient $-4$ est négatif, donc $g$ est croissante. Le modèle est $V(t)=1350\times\left(1+\dfrac{4}{100}\right)^t$. Après $6$ ans, $V(6)=1350\times\num{1,04}^6\approx\num{1708,18}$.\end{solution}

\mxexuse{7}

\begin{homework}Reprendre l'exercice d'entraînement 2 avec les données suivantes : remplacer la donnée initiale par $155$ et le taux ou la raison par $3$. Rédiger en quatre lignes : données, formule, calcul, conclusion.\end{homework}
\begin{solution}La correction dépend de l'exercice repris. Vérifier surtout la cohérence du modèle, le nombre d'itérations, l'arrondi et la phrase de conclusion.\end{solution}

\begin{exitticket}[3 minutes]Sur une demi-feuille : écrire la règle principale de la séance, traiter un exemple très court et indiquer son degré de confiance de 1 à 4. Question(s) essentielle(s) du jour : 1, 2 et 3.\end{exitticket}

\ifmxteacher
\begin{method}
\textbf{Conduite en 1 h :} 7 min d'automatismes ; 10 min de mise en commun et trace ; 35 min d'exercices essentiels ; 5 min de correction ; 3 min de billet de sortie.\\
\textbf{Conduite en 2 h :} ajouter 20 min de recherche de l'activité, 20 min d'institutionnalisation, un second exercice autonome et 10 min de correction comparée.\\
\textbf{Différenciation :} fournir aux élèves fragiles une fiche à quatre cases (données, modèle, calcul, conclusion). Demander aux élèves rapides de produire un contre-exemple, une justification ou une seconde méthode.
\end{method}
\fi
\mxsession{8}{\mxdate*{2026-10-08}}{B : 1 h -- C : 2 h}{Puissances et propriétés algébriques}{Fonctions exponentielles}
\begin{objectives}\begin{itemize}\item Transformer des puissances.\item Utiliser les exposants négatifs.\item Repérer une règle inapplicable.\end{itemize}\end{objectives}
\begin{fluency}[6 à 8 minutes, sans calculatrice]\begin{enumerate}\item Réduire $3x+4-x+2$. \mxblank{$2x+6$}\item Calculer $2^1\times2^3$. \mxblank{$2^4$}\item Écrire $5^{-2}$ sous forme de fraction. \mxblank{$\dfrac1{25}$}\item La fonction $x\mapsto\num{0,7}^x$ est-elle croissante ? \mxblank{Non, elle est décroissante.}\item Comparer $\num{1,04}^3$ et $\num{1,04}^7$. \mxblank{$\num{1,04}^3<\num{1,04}^7$}\end{enumerate}\end{fluency}
\begin{prerequisites}Relire la dernière trace écrite et savoir effectuer les calculs élémentaires mobilisés aujourd'hui. Avant l'activité, faire reformuler les mots-clés : transformer des puissances et utiliser les exposants négatifs.\end{prerequisites}

\begin{activity}
Simplifier $2^3\times2^{-5}$, $5^{x+2}\div5^x$ et $(3^x)^2$. Inventer ensuite un contre-exemple numérique à $a^x+a^y=a^{x+y}$.

\end{activity}
\begin{solution}On obtient $2^{-2}=\dfrac14$, $5^2=25$ et $3^{2x}$. Avec $a=2$ et $x=y=1$ on a $2+2=4$ : c'est une égalité accidentelle, qui ne prouve rien. Avec $a=3$ et $x=y=1$ : $3+3=6\neq9=3^{2}$, ce qui est un contre-exemple. Faire verbaliser la donnée qui déclenche le choix de méthode.\end{solution}

\begin{summary}$a^{x+y}=a^xa^y$, $a^{x-y}=\dfrac{a^x}{a^y}$ et $(a^x)^y=a^{xy}$. Ces règles exigent une même base positive.\end{summary}

\begin{example}
Reprendre les données de l'activité et rédiger la solution dans l'ordre : définition des objets, propriété utilisée, calcul exact ou approché, puis phrase de réponse. La réponse modèle figure dans l'édition enseignant.
\end{example}
\begin{solution}On obtient $2^{-2}=\dfrac14$, $5^2=25$ et $3^{2x}$. Avec $a=2$ et $x=y=1$ on a $2+2=4$ : c'est une égalité accidentelle, qui ne prouve rien. Avec $a=3$ et $x=y=1$ : $3+3=6\neq9=3^{2}$, ce qui est un contre-exemple.\end{solution}

\begin{exercise}[essentiel]
\begin{enumerate}\item Simplifier $3^{x+5}\div3^x$ et $(2^x)^3$.\item Étudier le sens de variation de $f(x)=5\times\num{1,05}^x$ et $g(x)=-5\times\num{0,8}^x$.\item Une valeur initiale de $1400$ augmente de $\pc{5}$ par an. Donner le modèle et calculer la valeur après $6$ ans.\end{enumerate}
\end{exercise}
\begin{solution}$3^5$ et $2^{3x}$. $f$ est croissante car $\num{1,05}>1$ et $5>0$ ; $\num{0,8}^x$ est décroissante et le coefficient $-5$ est négatif, donc $g$ est croissante. Le modèle est $V(t)=1400\times\left(1+\dfrac{5}{100}\right)^t$. Après $6$ ans, $V(6)=1400\times\num{1,05}^6\approx\num{1876,13}$.\end{solution}

\mxexuse{8}

\begin{homework}Reprendre l'exercice d'entraînement 3 avec les données suivantes : remplacer la donnée initiale par $160$ et le taux ou la raison par $4$. Rédiger en quatre lignes : données, formule, calcul, conclusion.\end{homework}
\begin{solution}La correction dépend de l'exercice repris. Vérifier surtout la cohérence du modèle, le nombre d'itérations, l'arrondi et la phrase de conclusion.\end{solution}

\begin{exitticket}[3 minutes]Sur une demi-feuille : écrire la règle principale de la séance, traiter un exemple très court et indiquer son degré de confiance de 1 à 4. Question(s) essentielle(s) du jour : 1 et 2.\end{exitticket}

\ifmxteacher
\begin{method}
\textbf{Conduite en 1 h :} 7 min d'automatismes ; 10 min de mise en commun et trace ; 35 min d'exercices essentiels ; 5 min de correction ; 3 min de billet de sortie.\\
\textbf{Conduite en 2 h :} ajouter 20 min de recherche de l'activité, 20 min d'institutionnalisation, un second exercice autonome et 10 min de correction comparée.\\
\textbf{Différenciation :} fournir aux élèves fragiles une fiche à quatre cases (données, modèle, calcul, conclusion). Demander aux élèves rapides de produire un contre-exemple, une justification ou une seconde méthode.
\end{method}
\fi
\mxsession{9}{\mxdate*{2026-10-13}}{B : 2 h -- C : 1 h}{Fonctions $x\mapsto k a^x$ : variations et graphique}{Fonctions exponentielles}
\begin{objectives}\begin{itemize}\item Étudier $ka^x$.\item Combiner signe et variation.\item Relier formule et graphique.\end{itemize}\end{objectives}
\begin{fluency}[6 à 8 minutes, sans calculatrice]\begin{enumerate}\item Calculer $\pc{20}$ de $100$. \mxblank{$20$}\item Écrire $5^{-2}$ sous forme de fraction. \mxblank{$\dfrac1{25}$}\item La fonction $x\mapsto\num{0,7}^x$ est-elle croissante ? \mxblank{Non, elle est décroissante.}\item Comparer $\num{1,04}^3$ et $\num{1,04}^7$. \mxblank{$\num{1,04}^3<\num{1,04}^7$}\item Calculer $2^2\times2^3$. \mxblank{$2^5$}\end{enumerate}\end{fluency}
\begin{prerequisites}Relire la dernière trace écrite et savoir effectuer les calculs élémentaires mobilisés aujourd'hui. Avant l'activité, faire reformuler les mots-clés : étudier $ka^x$ et combiner signe et variation.\end{prerequisites}

\begin{activity}
Les quatre courbes $\mathcal C_1$, $\mathcal C_2$, $\mathcal C_3$ et $\mathcal C_4$ ci-dessous représentent, dans un ordre inconnu, les fonctions $2^x$, $\num{0,6}^x$, $-2^x$ et $-\num{0,6}^x$. Associer chaque fonction à sa courbe. Justifier chaque choix par le signe, le point d'ordonnée à l'origine et le sens de variation.
\begin{center}\begin{tikzpicture}\begin{axis}[width=.68\textwidth,height=6.5cm,axis lines=middle,xmin=-3,xmax=3,ymin=-6,ymax=6,grid=major,xlabel={$x$},ylabel={$y$}]\addplot[mxteacher,very thick,samples=100,domain=-2.5:2.5]{2^x};\addplot[mxteacher,very thick,dashed,samples=100,domain=-2.5:2.5]{-2^x};\addplot[mxwarn,very thick,samples=100,domain=-2.5:2.5]{0.6^x};\addplot[mxwarn,very thick,dashed,samples=100,domain=-2.5:2.5]{-(0.6^x)};\node[mxwarn,anchor=south west] at (axis cs:-2.4,3.6) {$\mathcal C_1$};\node[mxteacher,anchor=north east] at (axis cs:1.5,-3.6) {$\mathcal C_2$};\node[mxteacher,anchor=south east] at (axis cs:2.3,4.9) {$\mathcal C_3$};\node[mxwarn,anchor=north west] at (axis cs:-2.4,-3.6) {$\mathcal C_4$};\end{axis}\end{tikzpicture}\end{center}
\end{activity}
\begin{solution}Toutes ont une valeur absolue égale à $1$ en $0$. Le signe distingue les deux paires ; la base et le signe du coefficient déterminent la variation. Association : $\mathcal C_3$ est $2^x$ (positive, croissante, passe par $\pt{0;1}$) ; $\mathcal C_1$ est $\num{0,6}^x$ (positive, décroissante, passe par $\pt{0;1}$) ; $\mathcal C_2$ est $-2^x$ (négative, décroissante, passe par $\pt{0;-1}$) ; $\mathcal C_4$ est $-\num{0,6}^x$ (négative, croissante, passe par $\pt{0;-1}$). Faire verbaliser la donnée qui déclenche le choix de méthode.\end{solution}

\begin{summary}Pour $f(x)=ka^x$, le signe est celui de $k$. Si $k<0$, la multiplication par $k$ inverse le sens de variation de $a^x$.\end{summary}

\begin{example}
Reprendre les données de l'activité et rédiger la solution dans l'ordre : définition des objets, propriété utilisée, calcul exact ou approché, puis phrase de réponse. La réponse modèle figure dans l'édition enseignant.
\end{example}
\begin{solution}Toutes ont une valeur absolue égale à $1$ en $0$. Le signe distingue les deux paires ; la base et le signe du coefficient déterminent la variation. Association : $\mathcal C_3$ est $2^x$ (positive, croissante, passe par $\pt{0;1}$) ; $\mathcal C_1$ est $\num{0,6}^x$ (positive, décroissante, passe par $\pt{0;1}$) ; $\mathcal C_2$ est $-2^x$ (négative, décroissante, passe par $\pt{0;-1}$) ; $\mathcal C_4$ est $-\num{0,6}^x$ (négative, croissante, passe par $\pt{0;-1}$).\end{solution}

\begin{exercise}[essentiel]
\begin{enumerate}\item Simplifier $3^{x+6}\div3^x$ et $(2^x)^3$.\item Étudier le sens de variation de $f(x)=6\times\num{1,06}^x$ et $g(x)=-6\times\num{0,8}^x$.\item Une valeur initiale de $1450$ augmente de $\pc{6}$ par an. Donner le modèle et calculer la valeur après $6$ ans.\end{enumerate}
\end{exercise}
\begin{solution}$3^6$ et $2^{3x}$. $f$ est croissante car $\num{1,06}>1$ et $6>0$ ; $\num{0,8}^x$ est décroissante et le coefficient $-6$ est négatif, donc $g$ est croissante. Le modèle est $V(t)=1450\times\left(1+\dfrac{6}{100}\right)^t$. Après $6$ ans, $V(6)=1450\times\num{1,06}^6\approx\num{2056,85}$.\end{solution}

\mxexuse{9}

\begin{homework}Reprendre l'exercice d'entraînement 2 avec les données suivantes : remplacer la donnée initiale par $165$ et le taux ou la raison par $5$. Rédiger en quatre lignes : données, formule, calcul, conclusion.\end{homework}
\begin{solution}La correction dépend de l'exercice repris. Vérifier surtout la cohérence du modèle, le nombre d'itérations, l'arrondi et la phrase de conclusion.\end{solution}

\begin{exitticket}[3 minutes]Sur une demi-feuille : écrire la règle principale de la séance, traiter un exemple très court et indiquer son degré de confiance de 1 à 4. Question(s) essentielle(s) du jour : 1, 2 et 3.\end{exitticket}

\ifmxteacher
\begin{method}
\textbf{Conduite en 1 h :} 7 min d'automatismes ; 10 min de mise en commun et trace ; 35 min d'exercices essentiels ; 5 min de correction ; 3 min de billet de sortie.\\
\textbf{Conduite en 2 h :} ajouter 20 min de recherche de l'activité, 20 min d'institutionnalisation, un second exercice autonome et 10 min de correction comparée.\\
\textbf{Différenciation :} fournir aux élèves fragiles une fiche à quatre cases (données, modèle, calcul, conclusion). Demander aux élèves rapides de produire un contre-exemple, une justification ou une seconde méthode.
\end{method}
\fi
\mxsession{10}{\mxdate*{2026-10-15}}{B : 1 h -- C : 2 h}{Taux d'évolution moyen}{Fonctions exponentielles}
\begin{objectives}\begin{itemize}\item Calculer un taux moyen.\item Distinguer taux global et moyen.\item Interpréter un exposant $\frac1n$.\end{itemize}\end{objectives}
\begin{fluency}[6 à 8 minutes, sans calculatrice]\begin{enumerate}\item Écrire le coefficient multiplicateur d'une hausse de $\pc{25}$. \mxblank{$1+\dfrac{25}{100}=\num{1,25}$}\item La fonction $x\mapsto\num{0,7}^x$ est-elle croissante ? \mxblank{Non, elle est décroissante.}\item Comparer $\num{1,04}^3$ et $\num{1,04}^7$. \mxblank{$\num{1,04}^3<\num{1,04}^7$}\item Calculer $2^3\times2^3$. \mxblank{$2^6$}\item Écrire $5^{-2}$ sous forme de fraction. \mxblank{$\dfrac1{25}$}\end{enumerate}\end{fluency}
\begin{prerequisites}Relire la dernière trace écrite et savoir effectuer les calculs élémentaires mobilisés aujourd'hui. Avant l'activité, faire reformuler les mots-clés : calculer un taux moyen et distinguer taux global et moyen.\end{prerequisites}

\begin{activity}
Le prix d'un service passe de $18$ à $\num{23,50}$ euros en cinq ans. Calculer le coefficient global, puis le coefficient annuel moyen et le taux annuel moyen. Vérifier en recomposant cinq fois l'évolution.

\end{activity}
\begin{solution}$Q=\dfrac{\num{23,5}}{18}\approx\num{1,3056}$, $q=Q^{\frac15}\approx\num{1,0548}$, soit environ $\pc{5,48}$ par an. Faire verbaliser la donnée qui déclenche le choix de méthode.\end{solution}

\begin{summary}Si $Q$ est le coefficient global de $n$ évolutions identiques, le coefficient moyen est $q=Q^{\frac1n}$ et le taux moyen est $q-1$.\end{summary}

\begin{example}
Reprendre les données de l'activité et rédiger la solution dans l'ordre : définition des objets, propriété utilisée, calcul exact ou approché, puis phrase de réponse. La réponse modèle figure dans l'édition enseignant.
\end{example}
\begin{solution}$Q=\dfrac{\num{23,5}}{18}\approx\num{1,3056}$, $q=Q^{\frac15}\approx\num{1,0548}$, soit environ $\pc{5,48}$ par an.\end{solution}

\begin{exercise}[essentiel]
\begin{enumerate}\item Simplifier $3^{x+2}\div3^x$ et $(2^x)^3$.\item Étudier le sens de variation de $f(x)=2\times\num{1,02}^x$ et $g(x)=-2\times\num{0,8}^x$.\item Une valeur initiale de $1500$ augmente de $\pc{2}$ par an. Donner le modèle et calculer la valeur après $6$ ans.\end{enumerate}
\end{exercise}
\begin{solution}$3^2$ et $2^{3x}$. $f$ est croissante car $\num{1,02}>1$ et $2>0$ ; $\num{0,8}^x$ est décroissante et le coefficient $-2$ est négatif, donc $g$ est croissante. Le modèle est $V(t)=1500\times\left(1+\dfrac{2}{100}\right)^t$. Après $6$ ans, $V(6)=1500\times\num{1,02}^6\approx\num{1689,24}$.\end{solution}

\mxexuse{10}

\begin{homework}Reprendre l'exercice d'entraînement 3 avec les données suivantes : remplacer la donnée initiale par $170$ et le taux ou la raison par $6$. Rédiger en quatre lignes : données, formule, calcul, conclusion.\end{homework}
\begin{solution}La correction dépend de l'exercice repris. Vérifier surtout la cohérence du modèle, le nombre d'itérations, l'arrondi et la phrase de conclusion.\end{solution}

\begin{exitticket}[3 minutes]Sur une demi-feuille : écrire la règle principale de la séance, traiter un exemple très court et indiquer son degré de confiance de 1 à 4. Question(s) essentielle(s) du jour : 1 et 2.\end{exitticket}

\ifmxteacher
\begin{method}
\textbf{Conduite en 1 h :} 7 min d'automatismes ; 10 min de mise en commun et trace ; 35 min d'exercices essentiels ; 5 min de correction ; 3 min de billet de sortie.\\
\textbf{Conduite en 2 h :} ajouter 20 min de recherche de l'activité, 20 min d'institutionnalisation, un second exercice autonome et 10 min de correction comparée.\\
\textbf{Différenciation :} fournir aux élèves fragiles une fiche à quatre cases (données, modèle, calcul, conclusion). Demander aux élèves rapides de produire un contre-exemple, une justification ou une seconde méthode.
\end{method}
\fi
\mxsession{11}{\mxdate*{2026-11-03}}{B : 2 h -- C : 1 h}{Lecture de représentations graphiques}{Fonctions exponentielles}
\begin{objectives}\begin{itemize}\item Lire images et antécédents.\item Comparer deux modèles.\item Estimer graphiquement.\end{itemize}\end{objectives}
\begin{fluency}[6 à 8 minutes, sans calculatrice]\begin{enumerate}\item Réduire $3x+7-x+2$. \mxblank{$2x+9$}\item Comparer $\num{1,04}^3$ et $\num{1,04}^7$. \mxblank{$\num{1,04}^3<\num{1,04}^7$}\item Calculer $2^4\times2^3$. \mxblank{$2^7$}\item Écrire $5^{-2}$ sous forme de fraction. \mxblank{$\dfrac1{25}$}\item La fonction $x\mapsto\num{0,7}^x$ est-elle croissante ? \mxblank{Non, elle est décroissante.}\end{enumerate}\end{fluency}
\begin{prerequisites}Relire la dernière trace écrite et savoir effectuer les calculs élémentaires mobilisés aujourd'hui. Avant l'activité, faire reformuler les mots-clés : lire images et antécédents et comparer deux modèles.\end{prerequisites}

\begin{activity}
Deux offres commerciales sont modélisées par $f(x)=120\times\num{1,08}^x$ et $g(x)=180\times\num{0,96}^x$. Tracer les courbes sur $[0,12]$, lire leur point d'intersection et interpréter les zones où chaque valeur est la plus grande.
\begin{center}\begin{tikzpicture}\begin{axis}[width=.72\textwidth,height=5.5cm,axis lines=middle,xmin=0,xmax=12,ymin=0,ymax=320,grid=major,xtick distance=2,ytick distance=50,xlabel={$x$},ylabel={$y$}]\addplot[mxteacher,very thick,samples=100,domain=0:12]{120*1.08^x};\addplot[mxwarn,very thick,samples=100,domain=0:12]{180*0.96^x};\node[mxteacher,anchor=south] at (axis cs:9.5,270) {$\mathcal C_f$};\node[mxwarn,anchor=south] at (axis cs:9,125) {$\mathcal C_g$};\end{axis}\end{tikzpicture}\end{center}
\end{activity}
\begin{solution}La lecture graphique situe l'intersection vers $x\approx\num{3,4}$ (valeur calculée : $\dfrac{\log\num{1,5}}{\log\num{1,125}}\approx\num{3,44}$), pour une ordonnée voisine de $\num{180}\times\num{0,96}^{3,44}\approx\num{156}$. Avant ce rang $g$ est plus grande que $f$ ; après, $f$ est plus grande que $g$. Une valeur approchée est attendue. Faire verbaliser la donnée qui déclenche le choix de méthode.\end{solution}

\begin{summary}Lire $f(x)$ consiste à partir de $x$ vers la courbe puis vers l'axe des ordonnées. Chercher un antécédent de $y$ consiste à partir de $y$ vers la courbe puis vers l'axe des abscisses.\end{summary}

\begin{example}
Reprendre les données de l'activité et rédiger la solution dans l'ordre : définition des objets, propriété utilisée, calcul exact ou approché, puis phrase de réponse. La réponse modèle figure dans l'édition enseignant.
\end{example}
\begin{solution}La lecture graphique situe l'intersection vers $x\approx\num{3,4}$ (valeur calculée : $\dfrac{\log\num{1,5}}{\log\num{1,125}}\approx\num{3,44}$), pour une ordonnée voisine de $\num{180}\times\num{0,96}^{3,44}\approx\num{156}$. Avant ce rang $g$ est plus grande que $f$ ; après, $f$ est plus grande que $g$. Une valeur approchée est attendue.\end{solution}

\begin{exercise}[essentiel]
\begin{enumerate}\item Simplifier $3^{x+3}\div3^x$ et $(2^x)^3$.\item Étudier le sens de variation de $f(x)=3\times\num{1,03}^x$ et $g(x)=-3\times\num{0,8}^x$.\item Une valeur initiale de $1550$ augmente de $\pc{3}$ par an. Donner le modèle et calculer la valeur après $6$ ans.\end{enumerate}
\end{exercise}
\begin{solution}$3^3$ et $2^{3x}$. $f$ est croissante car $\num{1,03}>1$ et $3>0$ ; $\num{0,8}^x$ est décroissante et le coefficient $-3$ est négatif, donc $g$ est croissante. Le modèle est $V(t)=1550\times\left(1+\dfrac{3}{100}\right)^t$. Après $6$ ans, $V(6)=1550\times\num{1,03}^6\approx\num{1850,78}$.\end{solution}

\mxexuse{11}

\begin{homework}Reprendre l'exercice d'entraînement 2 avec les données suivantes : remplacer la donnée initiale par $175$ et le taux ou la raison par $7$. Rédiger en quatre lignes : données, formule, calcul, conclusion.\end{homework}
\begin{solution}La correction dépend de l'exercice repris. Vérifier surtout la cohérence du modèle, le nombre d'itérations, l'arrondi et la phrase de conclusion.\end{solution}

\begin{exitticket}[3 minutes]Sur une demi-feuille : écrire la règle principale de la séance, traiter un exemple très court et indiquer son degré de confiance de 1 à 4. Question(s) essentielle(s) du jour : 1, 2 et 3.\end{exitticket}

\ifmxteacher
\begin{method}
\textbf{Conduite en 1 h :} 7 min d'automatismes ; 10 min de mise en commun et trace ; 35 min d'exercices essentiels ; 5 min de correction ; 3 min de billet de sortie.\\
\textbf{Conduite en 2 h :} ajouter 20 min de recherche de l'activité, 20 min d'institutionnalisation, un second exercice autonome et 10 min de correction comparée.\\
\textbf{Différenciation :} fournir aux élèves fragiles une fiche à quatre cases (données, modèle, calcul, conclusion). Demander aux élèves rapides de produire un contre-exemple, une justification ou une seconde méthode.
\end{method}
\fi
\mxsession{12}{\mxdate*{2026-11-05}}{B : 1 h -- C : 2 h}{Équations exponentielles simples}{Fonctions exponentielles}
\begin{objectives}\begin{itemize}\item Résoudre avec une même base.\item Approcher une solution.\item Vérifier la vraisemblance.\end{itemize}\end{objectives}
\begin{fluency}[6 à 8 minutes, sans calculatrice]\begin{enumerate}\item Calculer $\pc{5}$ de $160$. \mxblank{$8$}\item Calculer $2^1\times2^3$. \mxblank{$2^4$}\item Écrire $5^{-2}$ sous forme de fraction. \mxblank{$\dfrac1{25}$}\item La fonction $x\mapsto\num{0,7}^x$ est-elle croissante ? \mxblank{Non, elle est décroissante.}\item Comparer $\num{1,04}^3$ et $\num{1,04}^7$. \mxblank{$\num{1,04}^3<\num{1,04}^7$}\end{enumerate}\end{fluency}
\begin{prerequisites}Relire la dernière trace écrite et savoir effectuer les calculs élémentaires mobilisés aujourd'hui. Avant l'activité, faire reformuler les mots-clés : résoudre avec une même base et approcher une solution.\end{prerequisites}

\begin{activity}
Résoudre $4^x=64$, $9^x=\dfrac13$ et $\num{1,12}^x=2$ : d'abord par même base lorsque c'est possible, puis numériquement pour la dernière. Encadrer la dernière solution entre deux entiers.

\end{activity}
\begin{solution}$x=3$ ; $9^x=3^{-1}$ donc $2x=-1$ et $x=-\dfrac12$ ; enfin $x\approx\num{6,12}$, entre $6$ et $7$. Faire verbaliser la donnée qui déclenche le choix de méthode.\end{solution}

\begin{summary}Pour $a>0$ et $a\neq1$, pour résoudre $a^x=b$ : utiliser une même base quand c'est possible ; sinon encadrer ou utiliser un outil numérique. Vérifier que $b>0$.\end{summary}

\begin{example}
Reprendre les données de l'activité et rédiger la solution dans l'ordre : définition des objets, propriété utilisée, calcul exact ou approché, puis phrase de réponse. La réponse modèle figure dans l'édition enseignant.
\end{example}
\begin{solution}$x=3$ ; $9^x=3^{-1}$ donc $2x=-1$ et $x=-\dfrac12$ ; enfin $x\approx\num{6,12}$, entre $6$ et $7$.\end{solution}

\begin{exercise}[essentiel]
\begin{enumerate}\item Simplifier $3^{x+4}\div3^x$ et $(2^x)^3$.\item Étudier le sens de variation de $f(x)=4\times\num{1,04}^x$ et $g(x)=-4\times\num{0,8}^x$.\item Une valeur initiale de $1600$ augmente de $\pc{4}$ par an. Donner le modèle et calculer la valeur après $6$ ans.\end{enumerate}
\end{exercise}
\begin{solution}$3^4$ et $2^{3x}$. $f$ est croissante car $\num{1,04}>1$ et $4>0$ ; $\num{0,8}^x$ est décroissante et le coefficient $-4$ est négatif, donc $g$ est croissante. Le modèle est $V(t)=1600\times\left(1+\dfrac{4}{100}\right)^t$. Après $6$ ans, $V(6)=1600\times\num{1,04}^6\approx\num{2024,51}$.\end{solution}

\mxexuse{12}

\begin{homework}Reprendre l'exercice d'entraînement 3 avec les données suivantes : remplacer la donnée initiale par $180$ et le taux ou la raison par $2$. Rédiger en quatre lignes : données, formule, calcul, conclusion.\end{homework}
\begin{solution}La correction dépend de l'exercice repris. Vérifier surtout la cohérence du modèle, le nombre d'itérations, l'arrondi et la phrase de conclusion.\end{solution}

\begin{exitticket}[3 minutes]Sur une demi-feuille : écrire la règle principale de la séance, traiter un exemple très court et indiquer son degré de confiance de 1 à 4. Question(s) essentielle(s) du jour : 1 et 2.\end{exitticket}

\ifmxteacher
\begin{method}
\textbf{Conduite en 1 h :} 7 min d'automatismes ; 10 min de mise en commun et trace ; 35 min d'exercices essentiels ; 5 min de correction ; 3 min de billet de sortie.\\
\textbf{Conduite en 2 h :} ajouter 20 min de recherche de l'activité, 20 min d'institutionnalisation, un second exercice autonome et 10 min de correction comparée.\\
\textbf{Différenciation :} fournir aux élèves fragiles une fiche à quatre cases (données, modèle, calcul, conclusion). Demander aux élèves rapides de produire un contre-exemple, une justification ou une seconde méthode.
\end{method}
\fi
\mxsession{13}{\mxdate*{2026-11-10}}{B : 2 h -- C : 1 h}{Modélisation par une fonction exponentielle}{Fonctions exponentielles}
\begin{objectives}\begin{itemize}\item Construire un modèle exponentiel.\item Identifier les paramètres.\item Prévoir et interpréter.\end{itemize}\end{objectives}
\begin{fluency}[6 à 8 minutes, sans calculatrice]\begin{enumerate}\item Écrire le coefficient multiplicateur d'une hausse de $\pc{10}$. \mxblank{$1+\dfrac{10}{100}=\num{1,10}$}\item Écrire $5^{-2}$ sous forme de fraction. \mxblank{$\dfrac1{25}$}\item La fonction $x\mapsto\num{0,7}^x$ est-elle croissante ? \mxblank{Non, elle est décroissante.}\item Comparer $\num{1,04}^3$ et $\num{1,04}^7$. \mxblank{$\num{1,04}^3<\num{1,04}^7$}\item Calculer $2^2\times2^3$. \mxblank{$2^5$}\end{enumerate}\end{fluency}
\begin{prerequisites}Relire la dernière trace écrite et savoir effectuer les calculs élémentaires mobilisés aujourd'hui. Avant l'activité, faire reformuler les mots-clés : construire un modèle exponentiel et identifier les paramètres.\end{prerequisites}

\begin{activity}
Une application compte $\num{12000}$ utilisateurs au lancement et sa fréquentation augmente de $\pc{9}$ par mois. Construire $f(t)$, calculer $f(6)$, puis déterminer numériquement quand elle dépassera $\num{25000}$ utilisateurs.
\begin{center}\begin{tikzpicture}\begin{axis}[width=.72\textwidth,height=5.5cm,axis lines=middle,xmin=0,xmax=12,ymin=0,ymax=36000,grid=major,xtick distance=2,ytick distance=5000,scaled y ticks=false,xlabel={$t$ (mois)},ylabel={utilisateurs},ylabel style={at={(axis description cs:.02,.98)},anchor=south west}]\addplot[mxteacher,very thick,samples=100,domain=0:12]{12000*1.09^x};\addplot[mxwarn,dashed,thick,domain=0:12]{25000};\node[mxteacher,anchor=south] at (axis cs:9.2,31500) {$f$};\node[mxwarn,anchor=south west] at (axis cs:0.3,25500) {seuil $\num{25000}$};\end{axis}\end{tikzpicture}\end{center}
\end{activity}
\begin{solution}$f(t)=\num{12000}\times\num{1,09}^t$, $f(6)\approx\num{20125}$. Le seuil est dépassé dès $t=9$, c'est-à-dire $9$ mois après le lancement, car $f(8)\approx\num{23911}<\num{25000}<f(9)\approx\num{26063}$. Faire verbaliser la donnée qui déclenche le choix de méthode.\end{solution}

\begin{summary}Une évolution relative constante de coefficient $q$ se modélise par $f(t)=V_0q^t$. $V_0$ est la valeur initiale et $q$ le coefficient par unité de temps.\end{summary}

\begin{example}
Reprendre les données de l'activité et rédiger la solution dans l'ordre : définition des objets, propriété utilisée, calcul exact ou approché, puis phrase de réponse. La réponse modèle figure dans l'édition enseignant.
\end{example}
\begin{solution}$f(t)=\num{12000}\times\num{1,09}^t$, $f(6)\approx\num{20125}$. Le seuil est dépassé dès $t=9$, c'est-à-dire $9$ mois après le lancement, car $f(8)\approx\num{23911}<\num{25000}<f(9)\approx\num{26063}$.\end{solution}

\begin{exercise}[essentiel]
\begin{enumerate}\item Simplifier $3^{x+5}\div3^x$ et $(2^x)^3$.\item Étudier le sens de variation de $f(x)=5\times\num{1,05}^x$ et $g(x)=-5\times\num{0,8}^x$.\item Une valeur initiale de $1650$ augmente de $\pc{5}$ par an. Donner le modèle et calculer la valeur après $6$ ans.\end{enumerate}
\end{exercise}
\begin{solution}$3^5$ et $2^{3x}$. $f$ est croissante car $\num{1,05}>1$ et $5>0$ ; $\num{0,8}^x$ est décroissante et le coefficient $-5$ est négatif, donc $g$ est croissante. Le modèle est $V(t)=1650\times\left(1+\dfrac{5}{100}\right)^t$. Après $6$ ans, $V(6)=1650\times\num{1,05}^6\approx\num{2211,16}$.\end{solution}

\mxexuse{13}

\begin{homework}Reprendre l'exercice d'entraînement 2 avec les données suivantes : remplacer la donnée initiale par $185$ et le taux ou la raison par $3$. Rédiger en quatre lignes : données, formule, calcul, conclusion.\end{homework}
\begin{solution}La correction dépend de l'exercice repris. Vérifier surtout la cohérence du modèle, le nombre d'itérations, l'arrondi et la phrase de conclusion.\end{solution}

\begin{exitticket}[3 minutes]Sur une demi-feuille : écrire la règle principale de la séance, traiter un exemple très court et indiquer son degré de confiance de 1 à 4. Question(s) essentielle(s) du jour : 1, 2 et 3.\end{exitticket}

\ifmxteacher
\begin{method}
\textbf{Conduite en 1 h :} 7 min d'automatismes ; 10 min de mise en commun et trace ; 35 min d'exercices essentiels ; 5 min de correction ; 3 min de billet de sortie.\\
\textbf{Conduite en 2 h :} ajouter 20 min de recherche de l'activité, 20 min d'institutionnalisation, un second exercice autonome et 10 min de correction comparée.\\
\textbf{Différenciation :} fournir aux élèves fragiles une fiche à quatre cases (données, modèle, calcul, conclusion). Demander aux élèves rapides de produire un contre-exemple, une justification ou une seconde méthode.
\end{method}
\fi
\mxsession{14}{\mxdate*{2026-11-12}}{B : 1 h -- C : 2 h}{Bilan et évaluation du chapitre}{Fonctions exponentielles}
\begin{objectives}\begin{itemize}\item Mobiliser les acquis du chapitre.\item Rédiger une modélisation.\item Corriger une erreur.\end{itemize}\end{objectives}
\begin{fluency}[6 à 8 minutes, sans calculatrice]\begin{enumerate}\item Réduire $3x+3-x+2$. \mxblank{$2x+5$}\item La fonction $x\mapsto\num{0,7}^x$ est-elle croissante ? \mxblank{Non, elle est décroissante.}\item Comparer $\num{1,04}^3$ et $\num{1,04}^7$. \mxblank{$\num{1,04}^3<\num{1,04}^7$}\item Calculer $2^4\times2^3$. \mxblank{$2^7$}\item Écrire $5^{-2}$ sous forme de fraction. \mxblank{$\dfrac1{25}$}\end{enumerate}\end{fluency}
\begin{prerequisites}Relire la dernière trace écrite et savoir effectuer les calculs élémentaires mobilisés aujourd'hui. Avant l'activité, faire reformuler les mots-clés : mobiliser les acquis du chapitre et rédiger une modélisation.\end{prerequisites}

\begin{activity}
À partir de quatre productions d'élèves, repérer et corriger : confusion taux/coefficient, propriété de puissances fausse, erreur de sens de variation et conclusion sans unité.
\begin{enumerate}[label=\textbf{Production \arabic*.},leftmargin=*]
\item Une population de $\num{2000}$ habitants augmente de $\pc{3}$ par an. Léa écrit : « $P(n)=\num{2000}\times3^n$. »
\item Karim écrit : « $2^x+2^3=2^{x+3}$. »
\item Awa écrit : « $f(x)=-3\times\num{0,5}^x$ est décroissante, car $\num{0,5}<1$. »
\item Une machine achetée $\num{12000}$ euros perd $\pc{8}$ de sa valeur par an. Tojo écrit : « $V(n)=\num{12000}\times\num{0,92}^n$, donc $V(5)\approx\num{7909}$. Réponse : $\num{7909}$. »
\end{enumerate}

\end{activity}
\begin{solution}La correction doit nommer l'erreur, écrire la règle correcte, reprendre le calcul et conclure dans le contexte. Production 1 : confusion entre taux et coefficient ; une hausse de $\pc{3}$ correspond à $q=\num{1,03}$, donc $P(n)=\num{2000}\times\num{1,03}^n$. Production 2 : propriété de puissances fausse ; la règle porte sur les produits, $2^x\times2^3=2^{x+3}$, mais la somme $2^x+2^3$ ne se simplifie pas (avec $x=1$ : $2+8=10\neq16=2^4$). Production 3 : erreur de sens de variation ; $\num{0,5}^x$ est décroissante, mais le coefficient $-3$ est négatif et inverse le sens de variation, donc $f$ est croissante (par exemple $f(0)=-3<f(1)=-\num{1,5}$). Production 4 : conclusion sans unité ; le modèle et le calcul sont corrects, mais il faut conclure dans le contexte : après $5$ ans, la machine vaut environ $\num{7909}$ euros. Faire verbaliser la donnée qui déclenche le choix de méthode.\end{solution}

\begin{summary}Pour valider un modèle exponentiel : identifier $V_0$ et $q$, calculer avec les unités, vérifier l'ordre de grandeur et interpréter le résultat.\end{summary}

\begin{example}
Reprendre les données de l'activité et rédiger la solution dans l'ordre : définition des objets, propriété utilisée, calcul exact ou approché, puis phrase de réponse. La réponse modèle figure dans l'édition enseignant.
\end{example}
\begin{solution}La correction doit nommer l'erreur, écrire la règle correcte, reprendre le calcul et conclure dans le contexte. Production 1 : confusion entre taux et coefficient ; une hausse de $\pc{3}$ correspond à $q=\num{1,03}$, donc $P(n)=\num{2000}\times\num{1,03}^n$. Production 2 : propriété de puissances fausse ; la règle porte sur les produits, $2^x\times2^3=2^{x+3}$, mais la somme $2^x+2^3$ ne se simplifie pas (avec $x=1$ : $2+8=10\neq16=2^4$). Production 3 : erreur de sens de variation ; $\num{0,5}^x$ est décroissante, mais le coefficient $-3$ est négatif et inverse le sens de variation, donc $f$ est croissante (par exemple $f(0)=-3<f(1)=-\num{1,5}$). Production 4 : conclusion sans unité ; le modèle et le calcul sont corrects, mais il faut conclure dans le contexte : après $5$ ans, la machine vaut environ $\num{7909}$ euros.\end{solution}

\begin{exercise}[essentiel]
\begin{enumerate}\item Simplifier $3^{x+6}\div3^x$ et $(2^x)^3$.\item Étudier le sens de variation de $f(x)=6\times\num{1,06}^x$ et $g(x)=-6\times\num{0,8}^x$.\item Une valeur initiale de $1700$ augmente de $\pc{6}$ par an. Donner le modèle et calculer la valeur après $6$ ans.\end{enumerate}
\end{exercise}
\begin{solution}$3^6$ et $2^{3x}$. $f$ est croissante car $\num{1,06}>1$ et $6>0$ ; $\num{0,8}^x$ est décroissante et le coefficient $-6$ est négatif, donc $g$ est croissante. Le modèle est $V(t)=1700\times\left(1+\dfrac{6}{100}\right)^t$. Après $6$ ans, $V(6)=1700\times\num{1,06}^6\approx\num{2411,48}$.\end{solution}

\mxexuse{14}

\begin{homework}Reprendre l'exercice d'entraînement 3 avec les données suivantes : remplacer la donnée initiale par $190$ et le taux ou la raison par $4$. Rédiger en quatre lignes : données, formule, calcul, conclusion.\end{homework}
\begin{solution}La correction dépend de l'exercice repris. Vérifier surtout la cohérence du modèle, le nombre d'itérations, l'arrondi et la phrase de conclusion.\end{solution}

\begin{exitticket}[3 minutes]Sur une demi-feuille : écrire la règle principale de la séance, traiter un exemple très court et indiquer son degré de confiance de 1 à 4. Question(s) essentielle(s) du jour : 1 et 2.\end{exitticket}

\ifmxteacher
\begin{method}
\textbf{Conduite en 1 h :} 7 min d'automatismes ; 10 min de mise en commun et trace ; 35 min d'exercices essentiels ; 5 min de correction ; 3 min de billet de sortie.\\
\textbf{Conduite en 2 h :} ajouter 20 min de recherche de l'activité, 20 min d'institutionnalisation, un second exercice autonome et 10 min de correction comparée.\\
\textbf{Différenciation :} fournir aux élèves fragiles une fiche à quatre cases (données, modèle, calcul, conclusion). Demander aux élèves rapides de produire un contre-exemple, une justification ou une seconde méthode.
\end{method}
\fi
